\documentclass[aspectratio=169]{beamer}
\usetheme{Madrid}
\usecolortheme{seahorse}

\usepackage{graphicx}
\usepackage{booktabs}
\usepackage{amsmath}
\usepackage{amssymb}
\usepackage{tabularx}
\usepackage{array}
\usepackage{ragged2e}

\title[Chapter 4: Regression]{\textbf{Chapter 4: Regression}}
\subtitle{School of Mathematics and Computational Sciences \\ University of Prince Edward Island}
\author{Instructor: Ali Raisolsadat}
\date{Part I: Exploring Data}

\setbeamertemplate{navigation symbols}{}
\setbeamertemplate{itemize items}[circle]
\setbeamertemplate{enumerate items}[default]
\graphicspath{{./}{figures/}}

\newcommand{\lectureimage}[2][0.90\linewidth]{%
  \IfFileExists{#2}{%
    \includegraphics[width=#1]{#2}%
  }{%
    \fbox{\parbox[c][0.36\textheight][c]{0.88\linewidth}{%
      \centering\small
      Figure file: \texttt{\detokenize{#2}}\\[0.15cm]
      Generate this figure by knitting the revised Chapter 4 RMarkdown notes.
    }}%
  }%
}

\begin{document}

% Slide 1
\begin{frame}
  \titlepage
\end{frame}

% Slide 2
\begin{frame}{Why Regression Matters}
\begin{itemize}
	\item Chapter 3 introduced scatterplots and correlation as tools for describing the direction and strength of a linear association between two quantitative variables.
	\item Regression goes one step further. It gives an equation for a fitted line that can be used to describe how a response variable tends to change as an explanatory variable changes.
	\item Regression is important because it connects visual patterns, formulas, prediction and interpretation .
\end{itemize}
\end{frame}

% Slide 3
\begin{frame}{From Correlation to Regression}
Correlation and regression are related but they answer different questions.

\begin{table}
\centering
\small
\begin{tabularx}{\textwidth}{>{\RaggedRight\arraybackslash}p{0.24\textwidth} >{\RaggedRight\arraybackslash}p{0.32\textwidth} X}
\toprule
Idea & Correlation & Regression \\
\midrule
Main purpose & Describes the direction and strength of a linear association & Describes the fitted relationship and predicts the response \\
Roles of variables & The two variables are treated symmetrically & One variable is explanatory and one is the response \\
Units & Unitless & The slope has response units per explanatory-variable unit \\
Typical question & How strong is the linear association & What response is predicted for a given value of $x$ \\
\bottomrule
\end{tabularx}
\end{table}
\end{frame}

% Slide 4
\begin{frame}{The Roles of the Variables}
\begin{itemize}
	\item A regression analysis begins by assigning roles to the variables.
	\begin{itemize}
		  \item The explanatory variable $x$ is used to explain or predict.
 		 \item The response variable $y$ is the outcome being described or predicted.
	\end{itemize}
	\item These roles come from the scientific question rather than from the order of the columns in a dataset.
	\item A fitted line describes an average pattern. Two individuals with the same value of $x$ can still have different observed responses.
\end{itemize}
\end{frame}

% Slide 5
\begin{frame}{Regression as a Model}
\begin{itemize}
	\item A fitted regression line summarizes a linear pattern in the observed data.
	\item The fitted line does not guarantee that every point lies on the line. 
	\item It does not establish that changes in the explanatory variable cause changes in the response variable.
\end{itemize}
\end{frame}

% Slide 6
\begin{frame}{The Anatomy of a Straight Line}
\begin{itemize}
	\item A straight line relating $x$ to $y$ has the form
	$$y=a+bx$$
	\item The slope $b$ describes the change in $y$ when $x$ increases by one unit.
	\item The intercept $a$ is the value of $y$ when $x=0$.
\end{itemize}

\textbf{Example}
\begin{itemize}
	\item Consider the straight line
	$$y=3+2x$$
	\item For this line
	$$a=3 \qquad\text{and}\qquad b=2$$
	\item The line crosses the $y$-axis at $(0,3)$.
	\item Each one-unit increase in $x$ produces a two-unit increase in $y$.
	\begin{center}
	\begin{tabular}{c|cccc}
	$x$ & $0$ & $1$ & $2$ & $3$ \\
	\hline
	$y=3+2x$ & $3$ & $5$ & $7$ & $9$
	\end{tabular}
	\end{center}
\end{itemize}
\end{frame}

% Slide 7
\begin{frame}{How the Slope and Intercept Are Read}
For the straight line
$$y=a+bx$$
the intercept and slope describe different features of the line.

\begin{itemize}
	\item The \textbf{intercept} is $a$. Setting $x=0$ gives
	$$y=a+b(0)=a$$
	Therefore, the line crosses the $y$-axis at the point $(0,a)$.
	\item The \textbf{slope} is $b$. It can also be written as
	$$b=\frac{\text{change in }y}{\text{change in }x} =\frac{\Delta y}{\Delta x}$$
	\item If $b>0$, the line rises from left to right. If $b<0$, the line falls from left to right. If $b=0$, the line is horizontal.
\end{itemize}
\end{frame}

% Slide 8
\begin{frame}{From an Algebraic Line to a Regression Line}
A regression line uses the same algebraic form but the predicted response is written as $\hat y$.
\[
\hat y=a+bx
\]

The notation matters.
\begin{itemize}
  \item $y$ is an observed response taken from the data
  \item $\hat y$ is the response predicted by the fitted line
\end{itemize}

This distinction becomes very important when residuals are introduced.
\end{frame}

% Slide 9
\begin{frame}{Reading a Regression Equation}
\begin{block}{Question}
A greenhouse study uses the fitted equation
\[
\widehat{\text{height}}=4.5+1.8(\text{hours of light})
\]
to predict final plant height in centimetres from daily light exposure in hours.

What does the slope mean? What height does the model predict for a plant receiving 6 hours of light per day?
\end{block}
\end{frame}

% Slide 10
\begin{frame}{Additional Example: Solution}
\begin{itemize}
	\item The slope is 1.8 centimetres per hour of light.
	\item According to the fitted linear model, a one-hour increase in daily light exposure is associated with an increase of 1.8 centimetres in predicted final height.
	\item The intercept is 4.5 centimetres. It is the predicted height when the explanatory variable equals zero. Whether that interpretation is useful depends on the scientific setting.
	\item For 6 hours of light, the predicted height is
	$$\widehat{\text{height}}=4.5+1.8(6)$$
	$$\widehat{\text{height}}=15.3\text{ cm}$$
	\item The fitted model predicts a final height of 15.3 cm for a plant receiving 6 hours of light per day.
\end{itemize}
\end{frame}

% Slide 11
\section{Example 4.1: Predicting Fat Gain}
\begin{frame}{Example 4.1: Predicting Fat Gain}
\scriptsize
\begin{block}{Question}
\begin{itemize}
	\item A study followed 16 young adults who overate for eight weeks.
	\item The explanatory variable was the change in nonexercise activity (NEA) measured in Calories. The response variable was fat gain measured in kilograms.
	\item How can the fitted regression line be used to predict fat gain for a person whose NEA increases by 400 Calories?
\end{itemize}
\end{block}

\begin{columns}[T]
\begin{column}{0.48\textwidth}
\centering
\begin{tabular}{ccc}
\toprule
Volunteer & NEA Change & Fat Gain \\
 & (Calories) & (kg) \\
\midrule
1 & -94 & 4.2 \\
2 & -57 & 3.0 \\
3 & -29 & 3.7 \\
4 & 135 & 2.7 \\
5 & 143 & 3.2 \\
6 & 151 & 3.6 \\
7 & 245 & 2.4 \\
8 & 355 & 1.3 \\
\bottomrule
\end{tabular}
\end{column}
\begin{column}{0.48\textwidth}
\centering
\begin{tabular}{ccc}
\toprule
Volunteer & NEA Change & Fat Gain \\
 & (Calories) & (kg) \\
\midrule
9 & 392 & 3.8 \\
10 & 473 & 1.7 \\
11 & 486 & 1.6 \\
12 & 535 & 2.2 \\
13 & 571 & 1.0 \\
14 & 580 & 0.4 \\
15 & 620 & 2.3 \\
16 & 690 & 1.1 \\
\bottomrule
\end{tabular}
\end{column}
\end{columns}
\end{frame}

% Slide 12
\begin{frame}{Example 4.1: Explanation}
\begin{itemize}
	\item The fitted least-squares regression line for these data is
	$$\widehat{\text{fat gain}}=3.505-0.00344(\text{NEA change})$$
	\item A prediction at an NEA change of 400 Calories is obtained by substituting $x=400$ into the fitted equation.
	\item The value 400 lies inside the observed range of NEA changes, so this prediction is an interpolation rather than an extrapolation.
	\item The predicted fat gain is
	$$\widehat{\text{fat gain}}=3.505-0.00344(400)$$
	$$\widehat{\text{fat gain}}\approx2.13\text{ kg}$$
	\item The value 2.13 kg is the response predicted by the fitted line for a person whose NEA increases by 400 Calories.
\end{itemize}
\end{frame}

% Slide 13
\begin{frame}{Example 4.1: Scatterplot and Fitted Line}
  \centering
  \lectureimage[0.82\linewidth]{fig_ch4_regression.png}
\end{frame}

% Slide 14
\begin{frame}{Reading the Context}
\begin{itemize}
	\item The slope is $-0.00344$ kg per Calorie.
	\item Within the range of NEA values represented by the data, the fitted model predicts that fat gain decreases by about 0.00344 kg for each additional Calorie of NEA change.
	\item A 100-Calorie increase in NEA corresponds to a predicted decrease of about 0.344 kg in fat gain according to the fitted line.
	$$\widehat{\text{fat gain}}=3.505-0.00344(100)=3.505-0.344=3.161$$
\end{itemize}
\end{frame}

% Slide 15
\begin{frame}{The Meaning of the Intercept}
\begin{itemize}
	\item The intercept is 3.505 kg. It represents the predicted fat gain when NEA change is zero.
	$$\widehat{\text{fat gain}}=3.505\text{ kg when NEA change}=0$$
	\item In this example, the value $x=0$ lies inside the observed range, so the intercept has a meaningful interpretation.
	\item In other applications, the intercept may be mathematically necessary while having little practical meaning.
\end{itemize}
\end{frame}

% Slide 16
\begin{frame}{Least-Squares Regression}
\begin{itemize}
	\item A line drawn by eye is subjective. Different people could draw slightly different lines through the same scatterplot.
	\item Least-squares regression provides a precise mathematical rule for choosing one line.
	\item For an observation $(x_i,y_i)$, the fitted line gives a predicted response $\hat y_i$. The residual is the signed vertical difference between the observed response and the predicted response.
	$$e_i=y_i-\hat y_i$$
	\item A residual compares what was observed with what the fitted line predicted.
	\begin{itemize}
	  \item A positive residual means the observed point lies above the fitted line
	  \item A negative residual means the observed point lies below the fitted line
	  \item A residual close to zero means the fitted line predicted that observation well
	\end{itemize}
\end{itemize}
\end{frame}

% Slide 17
\begin{frame}{Residual Plot Illustration}
  \centering
  \lectureimage[0.82\linewidth]{fig_ch4_residuals.png}
\end{frame}

% Slide 18
\begin{frame}{Residuals from the NEA Study}
\begin{block}{Question}
Volunteer 8 had an NEA change of 355 Calories and an observed fat gain of 1.3 kg.
What fat gain does the fitted line predict for this volunteer and what is the residual?
\end{block}

Recall that $$\widehat{\text{fat gain}}=3.505-0.00344(\text{NEA change})$$
The fitted value is
$$\hat y=3.505-0.00344(355)$$
$$\hat y\approx2.28\text{ kg}$$

The residual is
$$e=y-\hat y=1.3-2.28$$
$$e\approx-0.98\text{ kg}$$

The residual is negative. So this point lies below the fitted line.
\end{frame}

% Slide 19
\begin{frame}{Why Are the Residuals Squared}
The least-squares line is the line that minimizes the sum of squared residuals.
$$\text{Sum of squared residuals}=\sum e_i^2$$

Squaring serves two purposes.
\begin{itemize}
  \item Positive and negative residuals do not cancel out
  \item Large prediction errors are penalized more strongly than small ones
\end{itemize}
\end{frame}

% Slide 20
\begin{frame}{Notation and Formula Reference}
\begin{block}{Least-Squares Regression Formulas}
$$b=r\frac{s_y}{s_x} \quad \text{(slope)}$$
$$a=\bar y-b\bar x  \quad \text{(intercept)}$$
$$\hat y=a+bx  \quad \text{(line equation)}$$
$$e=y-\hat y  \quad \text{(residuals)}$$
\end{block}

These formulas connect the regression coefficients to the summary statistics from Chapter 3.
\end{frame}

% Slide 21
\begin{frame}{What the Slope Formula Tells Us}
The slope formula
$$b=r\frac{s_y}{s_x}$$
shows that the slope depends on three ingredients.
\begin{itemize}
	\item The sign and strength of the linear association through $r$
	\item The spread of the response variable through $s_y$
	\item The spread of the explanatory variable through $s_x$
\end{itemize}
A positive correlation gives a positive slope. A negative correlation gives a negative slope.
\end{frame}

% Slide 22
\begin{frame}{Example 4.2: Calculating the Line from Summary Statistics}
\begin{block}{Question}
For the NEA and fat-gain data, the summary statistics are
$$\bar x=324.75$$
$$\bar y=2.3875$$
$$s_x=261.46$$
$$s_y=1.1207$$
$$r=-0.8021$$

Use these summary statistics to calculate the least-squares regression line.
\end{block}
\end{frame}

% Slide 23
\begin{frame}{Example 4.2: Calculate the Slope and Intercept}
\small
Using the slope formula
$$b=r\frac{s_y}{s_x}$$
$$b=(-0.8021)\left(\frac{1.1207}{261.46}\right)$$
$$b\approx-0.00344$$
The fitted line has a negative slope which agrees with the negative association seen in the scatterplot.

Using the intercept formula
$$a=\bar y-b\bar x$$
$$a=2.3875-(-0.00344)(324.75)$$
$$a\approx3.505$$

The least-squares regression line is therefore
$$\widehat{\text{fat gain}}=3.505-0.00344(\text{NEA change})$$
\end{frame}

% Slide 24
\begin{frame}{Example 4.2: Residual Calculations}
\scriptsize
A residual is the difference between the observed response and the fitted response:
$$e_i=y_i-\hat{y}_i$$

\begin{columns}[T,onlytextwidth]
\begin{column}{0.49\textwidth}
\centering
\setlength{\tabcolsep}{3.5pt}
\renewcommand{\arraystretch}{1.08}
\begin{tabular}{rrrrr}
\toprule
$i$ & $x_i$ & $y_i$ & $\hat{y}_i$ & $e_i$ \\
\midrule
$1$ & $-94$ & $4.2$ & $3.829$ & $0.371$ \\
$2$ & $-57$ & $3.0$ & $3.701$ & $-0.701$ \\
$3$ & $-29$ & $3.7$ & $3.605$ & $0.095$ \\
$4$ & $135$ & $2.7$ & $3.041$ & $-0.341$ \\
$5$ & $143$ & $3.2$ & $3.013$ & $0.187$ \\
$6$ & $151$ & $3.6$ & $2.985$ & $0.615$ \\
$7$ & $245$ & $2.4$ & $2.662$ & $-0.262$ \\
$8$ & $355$ & $1.3$ & $2.283$ & $-0.983$ \\
\bottomrule
\end{tabular}
\end{column}
\begin{column}{0.49\textwidth}
\centering
\setlength{\tabcolsep}{3.5pt}
\renewcommand{\arraystretch}{1.08}
\begin{tabular}{rrrrr}
\toprule
$i$ & $x_i$ & $y_i$ & $\hat{y}_i$ & $e_i$ \\
\midrule
$9$  & $392$ & $3.8$ & $2.156$ & $1.644$ \\
$10$ & $473$ & $1.7$ & $1.877$ & $-0.177$ \\
$11$ & $486$ & $1.6$ & $1.833$ & $-0.233$ \\
$12$ & $535$ & $2.2$ & $1.664$ & $0.536$ \\
$13$ & $571$ & $1.0$ & $1.540$ & $-0.540$ \\
$14$ & $580$ & $0.4$ & $1.509$ & $-1.109$ \\
$15$ & $620$ & $2.3$ & $1.371$ & $0.929$ \\
$16$ & $690$ & $1.1$ & $1.130$ & $-0.030$ \\
\bottomrule
\end{tabular}
\end{column}
\end{columns}

\vspace{0.2cm}

A positive residual means the observed value lies above the regression line. A negative residual means the observed value lies below the line. Using the unrounded fitted values
$$\sum_{i=1}^{16} e_i = 0.001$$
\end{frame}

% Slide 25
\begin{frame}[fragile]{Calculating the Least-Squares Line in R}
\scriptsize
\begin{verbatim}
# Enter the observed data
nea_change <- c( -94, -57, -29, 135, 143, 151, 245, 355, 392, 473, 486, 535, 571, 580, 620, 690)
fat_gain <- c(4.2, 3.0, 3.7, 2.7, 3.2, 3.6, 2.4, 1.3,3.8, 1.7, 1.6, 2.2, 1.0, 0.4, 2.3, 1.1)

# Calculate the required summary statistics
mean_x <- mean(nea_change)
mean_y <- mean(fat_gain)

sd_x <- sd(nea_change)
sd_y <- sd(fat_gain)

correlation <- cor(nea_change, fat_gain)

# Calculate the slope and intercept
slope <- correlation * (sd_y / sd_x)
intercept <- mean_y - slope * mean_x

> intercept     slope
>  3.505123 -0.003441
\end{verbatim}

Therefore, the least-squares regression line is
$\boxed{\hat{y}=3.5051-0.003441x}$
\end{frame}

% Slide 26
\begin{frame}[fragile]{R Code: Fitting the Least-Squares Line with \texttt{lm()}}
\scriptsize
R provides the \texttt{lm()} function for fitting a least-squares regression line.
\begin{block}{Fit the regression model}
\begin{verbatim}
# Fit fat gain as a function of NEA change
model <- lm(fat_gain ~ nea_change)

# Display the estimated coefficients
round(coef(model), digits = 6)

> (Intercept)  nea_change
>    3.505123   -0.003441
\end{verbatim}
\end{block}

\textbf{Reading an R model formula}
$$\texttt{response} \sim \texttt{explanatory variable}$$
\begin{itemize}
	\item The variable on the left of \texttt{\textasciitilde} is the response variable: \texttt{fat\_gain}.
	\item The variable on the right is the explanatory variable: \texttt{nea\_change}.
	\item The first coefficient is the intercept, and the second coefficient is the slope.
\end{itemize}

Therefore, R produces the regression line
$$\boxed{\widehat{\text{fat gain}}=3.5051-0.003441(\text{NEA change})}$$

The coefficients agree with the earlier hand calculation.
\end{frame}

% Slide 27
\begin{frame}[fragile]{R Code: Fitted Values and Residuals}
\scriptsize
\begin{verbatim}
# Fitted responses and residuals from the linear model
fitted_values <- fitted(model)
residual_values <- residuals(model)

# A standard summary of the fitted linear model
summary(model)
\end{verbatim}

At this stage, the fitted coefficients and $R^2$ are the most relevant parts of the output.
\end{frame}

% Slide 28
\begin{frame}{Break Time}
  \centering
  \lectureimage[0.35\linewidth]{meme_1.jpg}
\end{frame}

% Slide 29
\section{Facts About Least-Squares Regression}

\begin{frame}{Facts About Least-Squares Regression}
Several mathematical properties help explain how a fitted line behaves.
\begin{itemize}
  \item The roles of $x$ and $y$ are not interchangeable
  \item The sign of the slope always matches the sign of the correlation
  \item The least-squares line always passes through $(\bar x,\bar y)$
  \item The coefficient of determination $R^2$ summarizes how well the line accounts for variation in the response
\end{itemize}
\end{frame}

% Slide 30
\begin{frame}{Slope and Correlation}
The slope and correlation are connected by the formula
$$b=r\frac{s_y}{s_x}
$$
The sign of the slope always matches the sign of $r$.
\begin{itemize}
  \item If $r>0$, then $b>0$
  \item If $r<0$, then $b<0$
  \item If $r=0$, then $b=0$
\end{itemize}
\end{frame}

% Slide 31
\begin{frame}{The Roles of $x$ and $y$}
The least-squares line for predicting $y$ from $x$ is generally different from the least-squares line for predicting $x$ from $y$.

\vspace{0.15cm}
This difference occurs because the regression line minimizes vertical prediction errors in the response variable.

\vspace{0.15cm}
Regression is therefore directional in a way that correlation is not.
\end{frame}

% Slide 32
\begin{frame}{The Point of Averages}
The least-squares regression line always passes through the point of averages.
$$(\bar x,\bar y)$$

This fact is built into the intercept formula
$$a=\bar y-b\bar x$$

It provides a useful algebraic check when the fitted line is calculated by hand.
\end{frame}

% Slide 33
\begin{frame}{A Second View of $R^2$}
The variation in the observed response values can be divided into two parts: variation explained by the regression line and variation left in the residuals.
\begin{align*}
\text{SST} &=\sum_{i=1}^{n}(y_i-\bar{y})^2 &&\text{Total variation in the response}, \\[1mm]
\text{SSR} &=\sum_{i=1}^{n}(\hat{y}_i-\bar{y})^2 &&\text{Variation explained by the regression line}, \\[1mm]
\text{SSE} &=\sum_{i=1}^{n}(y_i-\hat{y}_i)^2 =\sum_{i=1}^{n}e_i^2 &&\text{Variation left unexplained}
\end{align*}
\end{frame}

% Slide 34
\begin{frame}{A Second View of $R^2$}
For a least-squares regression line with an intercept
$$\boxed{\text{SST}=\text{SSR}+\text{SSE}}$$

Therefore, the coefficient of determination can be written in two equivalent ways:
$$ \boxed{ R^2 =\frac{\text{SSR}}{\text{SST}} =1-\frac{\text{SSE}}{\text{SST}}} $$

\begin{itemize}
	\item $\dfrac{\text{SSR}}{\text{SST}}$ is the proportion of total variation explained by the regression line.
	\item $\dfrac{\text{SSE}}{\text{SST}}$ is the proportion of total variation left unexplained.
	\item For example, $R^2=0.64$ means that the linear model explains $64\%$ of the variation in the response variable. The remaining $36\%$ is left in the residuals.
\end{itemize}
\end{frame}

% Slide 35
\begin{frame}{Example 4.3: Interpreting $R^2$}
\begin{block}{Question}
Suppose a fitted linear model has
$$R^2=0.75$$

How should this value be interpreted?
\end{block}
\begin{itemize}
 	\item An $R^2$ value of 0.75 means that 75\% of the observed variation in the response variable is accounted for by the fitted linear relationship with the explanatory variable.
	\item The remaining 25\% is not accounted for by the fitted line and remains in the residual variation.
	\item This interpretation describes the quality of fit. It does not imply causation.
\end{itemize}
\end{frame}

% Slide 36
\begin{frame}{Outliers and Influential Observations}
A single unusual observation can affect correlation and least-squares regression in important ways.

\vspace{0.15cm}
Three related ideas should be distinguished.
\begin{itemize}
  \item An outlier has an unusually large residual
  \item A high-leverage point has an unusual value of $x$
  \item An influential point changes the fitted line noticeably when it is included or removed
\end{itemize}
\end{frame}

% Slide 37
\begin{frame}{Examples 4.4 and 4.5: The Empathy Study}
\begin{block}{Question}
An empathy study measures brain activity and empathy score for 16 women. Subject 16 has an unusually large empathy score.

How does this point affect the correlation and the least-squares regression line?
\end{block}

\begin{itemize}
	\item Subject 16 changes the correlation substantially.
	$$r=0.515\text{ with Subject 16}$$
	$$r=0.331\text{ without Subject 16}$$
	\item The fitted regression line also changes, although less dramatically than the correlation.
	\item This example shows that a point can be influential for both correlation and regression, with different levels of impact.
\end{itemize}
\end{frame}

% Slide 38
\begin{frame}{Empathy Study Plot}
  \centering
  \lectureimage[0.83\linewidth]{fig_ch4_outliers.png}
\end{frame}

% Slide 39
\begin{frame}{A Hypothetical High-Leverage Point}
A point with an extreme $x$ value has leverage because it can pull the fitted line toward itself.

\vspace{0.15cm}
If a high-leverage point also has a large residual, it can become highly influential.

\vspace{0.15cm}
This is why unusual observations should be examined carefully before a fitted line is interpreted.
\end{frame}

%% ======================
%% Slide 52
%% ======================
%\begin{frame}{Why Use a Logarithm Transformation}
%In biology and other sciences, data often span several orders of magnitude.
%
%\vspace{0.15cm}
%A logarithm transformation can help when
%\begin{itemize}
%  \item the raw data are strongly skewed
%  \item a few very large observations dominate the scale
%  \item a multiplicative relationship becomes more nearly linear after transformation
%\end{itemize}
%\end{frame}

%% ======================
%% Slide 53
%% ======================
%\begin{frame}{A Simple Base-10 Log Illustration}
%A base-10 logarithm compresses large scales.
%\[
%	\log_{10}(10)=1
%\]
%\[
%	\log_{10}(100)=2
%\]
%\[
%	\log_{10}(1000)=3
%\]
%\[
%	\log_{10}(10000)=4
%\]
%
%Values that are very far apart on the original scale become much closer on the log scale.
%\end{frame}

%% ======================
%% Slide 54
%% ======================
%\begin{frame}{Example 4.6: Mammal Brains and Body Weight}
%\begin{block}{Question}
%How does an animal's brain size relate to its body size across many mammal species?
%
%Why might a log-log transformation give a clearer regression picture than the raw variables?
%\end{block}
%\end{frame}

%% ======================
%% Slide 55
%% ======================
%\begin{frame}{Example 4.6: Explanation}
%On the raw scale, a few very large mammals dominate the axes and compress the smaller species near the origin.
%
%\vspace{0.15cm}
%After taking base-10 logarithms of both variables, the observations are spread more evenly and the relationship becomes more nearly linear.
%
%\vspace{0.15cm}
%The fitted log-log model is
%\[
%\log_{10}(\text{brain weight})=1.01+0.72\log_{10}(\text{body weight})
%\]
%\end{frame}

%% ======================
%% Slide 56
%% ======================
%\begin{frame}{Mammal Example Plot}
%  \centering
%  \lectureimage[0.88\linewidth]{fig_ch4_log_comparison.png}
%\end{frame}

%% ======================
%% Slide 57
%% ======================
%\begin{frame}{Back to the Original Scale}
%The log-log model can be rewritten as a power relationship.
%\[
%	y=10^{1.01}x^{0.72}
%\]
%
%This form shows that a multiplicative relationship on the original scale becomes linear after taking logarithms.
%\end{frame}

% Slide 40
\begin{frame}{Cautions About Correlation and Regression}
A fitted line should always be interpreted in light of the data and the scientific question.
\begin{itemize}
  \item Correlation and simple linear regression describe linear patterns
  \item Unusual observations can change the correlation and the fitted line
  \item Predictions outside the observed range require caution
  \item A lurking variable can help explain an observed association
  \item Association by itself does not imply causation
\end{itemize}
\end{frame}

% Slide 41
\begin{frame}{Linear Relationships}
A scatterplot should be examined before a simple linear model is interpreted.

\vspace{0.15cm}
Curvature, separate clusters or changing spread can all indicate that a straight line is an inadequate summary.

\vspace{0.15cm}
R or other statistical software can always produce a correlation coefficient and a fitted line, but those outputs are only useful when the relationship is reasonably linear.
\end{frame}

% Slide 42
\begin{frame}{Interpolation and Extrapolation}
A prediction within the observed range of $x$ values is called interpolation.

\vspace{0.15cm}
A prediction outside the observed range is called extrapolation.

\vspace{0.15cm}
Algebra allows the fitted line to be extended indefinitely, but the data do not show whether the same relationship continues beyond the observed range.
\end{frame}

% Slide 43
\begin{frame}{Additional Example: Interpolation versus Extrapolation}
\begin{block}{Question}
A fertilizer study observes daily doses from 0 to 5 millilitres and fits a linear regression model over that range.

Which prediction is better supported by the observed data: a prediction at 3 millilitres or a prediction at 100 millilitres?
\end{block}
\begin{itemize}
	\item The prediction at 3 millilitres is better supported because it is an interpolation.
	\item The prediction at 100 millilitres is an extrapolation far beyond the observed data.
	\item The fitted line contains no information about whether the biological relationship remains linear from 5 to 100 millilitres.
\end{itemize}
\end{frame}

% Slide 44
\begin{frame}{Lurking Variables}
A lurking variable is a variable that is not included as an explanatory or response variable in the analysis but may help explain the observed association.

\vspace{0.15cm}
For example, an observational study might find that coffee consumption is positively associated with risk of heart attack. In this case, smoking could be a lurking variable if it is related to both coffee consumption and heart attack risk.
\end{frame}

% Slide 45
\begin{frame}{Association Does Not Imply Causation}
\begin{itemize}
	\item Regression and correlation describe associations. They do not by themselves establish cause and effect.
	\item A causal conclusion requires more than a strong correlation, a steep slope or a large $R^2$. The study design matters.
	\item A well designed randomized experiment can provide stronger evidence for a causal effect because random assignment helps balance lurking variables across treatment groups.
	\item Observational studies generally require more caution because unmeasured confounding can remain even after statistical adjustment.
\end{itemize}
\end{frame}

% Slide 46
\begin{frame}{A Consistent Regression Workflow}
A useful first-year workflow is to interpret regression in the same order every time.
\begin{enumerate}
  \item Identify the explanatory variable and the response variable
  \item Examine the scatterplot for linearity and unusual observations
  \item Write and interpret the fitted line
  \item Calculate or interpret a prediction when appropriate
  \item Use residuals and $R^2$ to assess the fit
  \item Consider limitations such as extrapolation, lurking variables and causation
\end{enumerate}
\end{frame}

\end{document}
